\documentclass[12pt]{article}
\usepackage[utf8]{inputenc}
\usepackage[T1]{fontenc}
\usepackage[margin=1in]{geometry}
\usepackage{setspace}
\usepackage{xcolor}
\usepackage[colorlinks=true, urlcolor=blue!60!black, linkcolor=blue!60!black]{hyperref}
\setlength{\parindent}{2em}
\setlength{\parskip}{1.25em}
\renewcommand{\baselinestretch}{1.0}
\begin{document}
\begin{center}
{\LARGE \textit{Garcetti v.\ Ceballos}, 547 U.S.\ 410 (2006)\par}
\vspace{0.8em}
{\large Case Brief --- Four-Page Edition\par}
\end{center}
\vspace{1em}

\section{Memory Jogger}

A deputy district attorney who wrote an internal memo questioning a search-warrant affidavit spoke pursuant to his official duties, not as a citizen, so the First Amendment did not protect him from employer discipline (\href{https://supreme.justia.com/cases/federal/us/547/410/}{Justia}).

\section{Detailed Case Facts}

Richard Ceballos was a calendar deputy, with supervisory duties, in the Pomona branch of the Los Angeles County District Attorney's Office. In February 2000 a defense attorney told him that the affidavit supporting a critical search warrant was inaccurate. Ceballos examined the affidavit and visited the site, concluding that it called a separate roadway a ``long driveway'' and described tire tracks that the roadway's composition made difficult or impossible. On March 2, 2000, he submitted a disposition memorandum to his supervisor recommending dismissal. A heated meeting followed in which a sheriff's lieutenant sharply criticized him; the office proceeded with the prosecution, and the defense called Ceballos to testify at the suppression hearing. Ceballos alleged that District Attorney Gil Garcetti and his supervisors then retaliated by reassigning him to trial deputy, transferring him to the El Monte branch, and denying him a promotion. He sued under 42 U.S.C.\ \S~1983.

\section{Procedural History}

The district court granted summary judgment for the supervisors, holding the memo unprotected because it was written pursuant to Ceballos's job duties. The Ninth Circuit reversed under \textit{Pickering} and \textit{Connick}, finding the memo addressed a matter of public concern; Judge O'Scannlain specially concurred, urging no protection for routine job-duty speech. After reargument following Justice O'Connor's retirement, the Supreme Court reversed and remanded.

\section{Judicial Votes}

\textbf{5--4}, reversed and remanded (May 30, 2006).
\begin{itemize}
\item \textbf{Majority:} Justice Kennedy, joined by Chief Justice Roberts and Justices Scalia, Thomas, and Alito.
\item \textbf{Dissents:} Justice Stevens; Justice Souter, joined by Justices Stevens and Ginsburg; Justice Breyer.
\end{itemize}

\section{Holding}

When public employees make statements pursuant to their official duties, they do not speak as citizens, and the First Amendment does not insulate their communications from employer discipline. This threshold precedes \textit{Pickering} balancing, which applies only to speech as a citizen on a matter of public concern. Because the parties agreed the memo was part of Ceballos's job, the Court did not articulate a comprehensive test, but said the inquiry is ``a practical one'' and reserved whether the rule applies to ``speech related to scholarship or teaching.''

\section{Analysis of Opinions}

\textbf{Majority (Kennedy).} Neither the internal setting of the speech nor its job-related subject was dispositive; the controlling factor was that the memo was made pursuant to official duties. Drawing on \textit{Rosenberger}, the Court treated such speech as work product the employer ``commissioned or created.'' Official communications have official consequences, and constitutionalizing them would invite permanent judicial supervision of agency operations. Whistleblower laws, labor codes, ethics rules, and \textit{Brady} obligations provide alternative safeguards.

\textbf{Stevens (dissenting).} The right answer is ``Sometimes,'' not ``Never.'' Making protection turn on a job description is ``senseless'' and perversely encourages employees to go to the press rather than to their superiors.

\textbf{Souter, joined by Stevens and Ginsburg (dissenting).} Official-duty speech should be protected when it addresses matters of ``unusual importance,'' such as official dishonesty or threats to health and safety. The majority's ``bloated notion of controllable government speech'' could imperil academic freedom, and whistleblower statutes are an inconsistent ``patchwork.''

\textbf{Breyer (dissenting).} \textit{Pickering} balancing should apply only where a heightened need for protection meets diminished risk of judicial interference---as here, given lawyers' professional canons and \textit{Brady} duties.

\section{Examples: Future Applications}

\begin{enumerate}
\item \textbf{Same-Side: Building Inspector.} An inspector demoted for an internal report on a developer's fire-code violations has no First Amendment claim; filing such reports is her job.
\item \textbf{Same-Side: K--12 Teacher.} A teacher suspended for a personal political critique during an assigned class is generally treated as engaged in unprotected official-duty speech.
\item \textbf{Opposite-Side: Sworn Testimony.} A college official fired for subpoenaed trial testimony about fraud he found on the job engaged in citizen speech (\textit{Lane v.\ Franks}, 2014).
\item \textbf{Opposite-Side: Coach's Prayer.} A coach's quiet personal prayer after games was not official-duty speech (\textit{Kennedy v.\ Bremerton}, 2022).
\item \textbf{Fence-Sitter: Professor's Scholarship.} A public university professor denied a raise for a peer-reviewed article falls within the reserved scholarship question; the result depends on whether the circuit recognizes an academic-freedom exception.
\end{enumerate}

\section{Critique}

\textbf{Structure versus speech.} Formalist supporters say the First Amendment restrains the government as sovereign, not as employer, and that a bright line keeps federal judges from micromanaging agencies. Other critics say \textit{Garcetti} overextends government-speech doctrine and that its ``practical'' inquiry has produced circuit splits.

\textbf{Perverse incentives.} Progressive scholars such as Helen Norton argue the state may claim employees' speech as its own without accountability, depriving the public of insider information. Employees who report up the chain lose protection while leakers may keep it, and statutory whistleblower remedies are riddled with gaps.

\section{Key Quotations}

\begin{quote}
``We hold that when public employees make statements pursuant to their official duties, the employees are not speaking as citizens for First Amendment purposes, and the Constitution does not insulate their communications from employer discipline.''
\end{quote}
\noindent --- Justice Kennedy, for the Court, 547 U.S.\ at 421.

\begin{quote}
``Restricting speech that owes its existence to a public employee's professional responsibilities does not infringe any liberties the employee might have enjoyed as a private citizen. It simply reflects the exercise of employer control over what the employer itself has commissioned or created.''
\end{quote}
\noindent --- Justice Kennedy, for the Court, 547 U.S.\ at 421--422.

\begin{quote}
``The proper inquiry is a practical one. Formal job descriptions often bear little resemblance to the duties an employee actually is expected to perform, and the listing of a given task in an employee's written job description is neither necessary nor sufficient to demonstrate that conducting the task is within the scope of the employee's professional duties for First Amendment purposes.''
\end{quote}
\noindent --- Justice Kennedy, for the Court, 547 U.S.\ at 424--425.

\end{document}
